\chapter*{국문초록}

\begin{center}
  \begin{minipage}{0.8\textwidth}
    본 연구는 St. Petersburg 역설을 해결하고 인간의 의사결정 과정을 분석하기 위해 시행 횟수-이익 확률 모델을 제안한다. 기존의 효용 함수나 외부적 제약에 의존한 접근과 달리, 본 연구는 반복 시행에서의 이익 확률에 초점을 맞춘다. 연구 결과, 이 역설은 무한한 기대값에서 비롯된 것이 아니라 희박한 보상 구조에서 이익 확률의 수렴 속도가 느리다는 데 기인함을 밝혔다. 참가 비용이 높아질수록 이익 확률의 수렴 속도는 더욱 감소하여, 높은 비용의 게임을 비합리적으로 판단하는 이유를 설명할 수 있었다. 본 모델은 낮은 확률과 높은 보상이 결합된 상황에서의 의사결정을 해석할 수 있는 새로운 틀을 제공하며, St. Petersburg 역설에 대한 통찰을 제시한다.
  \end{minipage}
\end{center}

\vspace{1em}
주요어 : 상트페테르부르크의 역설, 의사결정 이론, 기대효용이론
